\section{The big picture: what question does this repo answer, and why care?}
\label{sec:bigpicture}

\subsection{The scientific question in one breath}
\label{sec:question}

How well does the loose, dusty soil on the Moon (called \textbf{regolith})
conduct heat, about a metre below the surface, \emph{at two specific places} ---
the Apollo~15 and Apollo~17 landing sites? And is that number \emph{the same at
both sites}, or different?

The single number that captures ``how well does deep Moon soil conduct heat'' is
called \Kd{} --- the \emph{deep thermal conductivity}. This whole repository
exists to measure \Kd{} at those two sites and to ask whether they differ.

\subsection{Why this is hard, and why it matters}

\begin{itemize}[leftmargin=1.4em,itemsep=3pt]
  \item The Moon has been mapped from orbit in incredible detail. But orbiters
    only see the \textbf{surface}. To know what is happening a metre \emph{down},
    you need a thermometer \emph{in the ground}.
  \item Thermometers have been buried in the Moon \textbf{exactly twice in
    history}: during Apollo~15 (1971) and Apollo~17 (1972). Astronauts drilled
    holes and dropped temperature probes down them. This was the
    \textbf{Heat-Flow Experiment (HFE)}. The instruments radioed back
    temperatures from 1971 to 1977.
  \item That 1971--1977 record is therefore \emph{the only direct measurement of
    metre-scale lunar soil conductivity that exists}. Everything else is an
    educated guess extrapolated from surface data.
  \item Current models (notably the widely used \textbf{Hayne et al.\ 2017}
    model) apply a \textbf{single conductivity value to the entire Moon}. This
    project asks: is one global number actually good enough, or do different
    places genuinely differ?
\end{itemize}

\begin{plainenglish}
We have two underground thermometers' worth of data, ever, from the whole Moon.
This project squeezes that scarce data as hard as it can to answer one sharp
question: does deep Moon soil conduct heat the same everywhere, or not?
\end{plainenglish}

\paragraph{Why anyone cares.}
\begin{itemize}[leftmargin=1.4em,itemsep=3pt]
  \item Future missions want to detect \textbf{water ice} trapped below the
    lunar surface. Whether ice stays frozen or sublimates away depends on the
    \textbf{temperature profile underground}, which depends directly on \Kd{}.
  \item An upcoming terahertz mission (\textbf{TSUKIMI}) will try to ``see'' into
    the regolith from orbit. To interpret what it sees, it needs the underground
    temperature-versus-depth curve $T(z)$ as a boundary condition --- and that
    curve is set by \Kd{}.
  \item More fundamentally: if ``one number for the whole Moon'' is wrong, every
    model built on it inherits the error.
\end{itemize}

\subsection{The headline result (the answer)}
\label{sec:headline}

After all the modelling and statistics, this is what the repository finds (these
exact numbers are committed in \file{output/kd\_retrieval\_results.json} and
quoted in the paper abstract).

\begin{table}[h]
\centering
\caption{The retrieved deep conductivity at each Apollo site.}
\label{tab:headline}
\begin{tabular}{lcc}
\toprule
\textbf{Site} & \textbf{Retrieved \Kdstar} & \textbf{Units} \\
\midrule
Apollo 15 & \textbf{4.58} \;(\textbf{+1.33} / $-$0.28) & \mwmk \\
Apollo 17 & \textbf{8.12} \;(\textbf{+0.49} / $-$0.61) & \mwmk \\
\bottomrule
\end{tabular}
\end{table}

\begin{itemize}[leftmargin=1.4em,itemsep=3pt]
  \item The little $*$ in \Kdstar{} means ``the best-fit value'' (the star is the
    retrieved estimate).
  \item The units \mwmk{} are milliwatts per metre per kelvin --- a measure of
    conductivity (\autoref{sec:conductivity} explains what conductivity units
    mean). For scale, the ``one global value'' everyone used before is
    \Kd{}~$=$~\num{3.4} in the same units, window glass is ${\sim}1000$, and
    copper is ${\sim}\num{4e8}$. Lunar regolith is an \emph{extraordinary}
    insulator.
  \item The \textbf{inter-site contrast} (the difference, A17~$-$~A15) is
    ${\approx}\,$\num{3.3}~\mwmk{}, with a 95\% confidence interval of roughly
    $[0.4, 4.6]$ that \textbf{excludes zero} ($p \approx 0.01$). In plain terms:
    the two sites really do appear to differ --- Apollo~17 soil conducts heat
    noticeably better than Apollo~15 soil --- by roughly a factor of~1.8.
\end{itemize}

\begin{keyidea}
The two sites are different. Apollo~17 deep soil conducts heat about
\textbf{1.8 times better} than Apollo~15 deep soil. That difference, established
from the only underground data that exist, is the scientific contribution of the
whole project. The rest of this guide explains \emph{how} the repository gets
there.
\end{keyidea}
